%% ============================================================================
%% sections/08-limitations.tex -- Limitations
%% Class-independent section file: it is \input by main-acm.tex and main-ieee.tex.
%% Rules (see README-build.md): no class-specific commands; cite with \cite only;
%% figure sizes come from the float budget and table bodies from tables/, never from this file.
%% Every paragraph and every caption is one numbered unit that starts with a
%% marker line "%% P:<id>".  Ids are frozen: never renumber, never reuse.
%% The same ids appear in paper.zh-CN.md; check_align.py compares the two.
%% Word budget of this section: 160 words of running text.
%% ============================================================================

\section{Limitations}\label{sec:limits}

%% P:sec8-p1
\ptag{sec8-p1}%
Physical results are \evPostRoute{}s on the ASAP7 predictive PDK at one process corner, with generated SRAM timing models; we claim no silicon, sign-off or frequency.
Each result names its \Lii storage contract; we never pair performance under the behavioral dual-port model with area and timing under single-port macros.
The text describes the target design; \secref{sec:uarch} states what the prototype lacks.

%% P:sec8-p2
\ptag{sec8-p2}%
System-level results are model projections; calibration on one PE does not transfer to the network, shared storage or off-chip memory.
Workload numbers are analytical counts of one decode step of one model family, not measurements.
Numerical experiments use synthetic tensors: they report numerical error under stated distributions and say nothing about model quality.
The same holds where the target design departs from the reference implementation: it quantizes the attention-core query to 8 bits (reference: BF16), and its online softmax and top-$k$ selection may differ from the reference kernels.
Fusion is evaluated only with a scalar epilogue, not in the block domain.
The third consumer tier is a minimal instance, not a complete RMSNorm.
